% =====================================================================
%  Chapter 23 — The Strategic Allocation Is a Decision, Not an Estimate
% =====================================================================

\chapter{The Strategic Allocation Is a Decision, Not an Estimate}\idx{Strategic allocation}
\label{ch:strategic}

The book's architecture has one rule that everything else follows from:
the strategic allocation is never set by the model. This chapter states
the rule, argues for it, and describes what the model is permitted to do
instead.

\section{The rule}

The system has four sleeves. The strategic allocation --- the split of the
portfolio across those four sleeves --- is fixed in advance by a stated
policy:

\begin{equation}
(w_{\mathrm{US}}, w_{\mathrm{intl}}, w_{\mathrm{bonds}}, w_{\mathrm{real}})
= (0.36,\ 0.24,\ 0.30,\ 0.10).
\end{equation}

The model cannot change these numbers. It cannot change them by a
percentage point. The only thing that varies them is the availability of
the underlying funds: a sleeve\idx{Sleeve} with no listed fund has its weight
redistributed over the sleeves that do have one, which is a fact about
the data rather than a decision.

What the model \emph{is} permitted to do is reallocate \emph{within} a
sleeve, subject to a bound. The bound is the fraction of the sleeve that
the tilt may move: at the book's default, a tilt can move a sleeve's
weight by up to an eighth in either direction, and the sleeve's total is
then restored.

\section{Why the rule is the right one}

There are three reasons, and the third is the important one.

The first is that the strategic allocation is a statement about the
investor's circumstances, not about the market. How much risk a person
should hold depends on their income, their age, the stability of their
employment and the horizon over which they will not need the money. None
of those are in the price series, and a model that could read them would
be a model of something other than returns.

The second is that the strategic allocation is the part of the portfolio
that the sequence-of-returns argument of Chapter 7 acts upon. A saver's
ruin probability is dominated by the allocation, not by the choice of
funds within a sleeve, and a system that lets a forecast choose the
allocation would be letting a quantity with no representation of the
sequence choose the thing that the sequence most affects.

The third is the one that decides it: \textbf{the strategic allocation is
the part of the portfolio that can be justified without a backtest}. The
reader can read the argument in Chapter 24 and disagree with it. The
reader cannot read the argument for the tilt, because there is none --- the
tilt is justified by a forecast whose value the reader cannot assess in
advance and whose record the book will show to be worth approximately
nothing.

\begin{takeaway}[The rule]
The strategic allocation is stated, argued, and fixed. The model may tilt
within a sleeve, by a bounded amount, and may not touch the split across
sleeves at all. Everything the model does is therefore reversible by the
reader without changing the book's argument.
\end{takeaway}

\section{What the model is permitted to do, precisely}

Three things, and no more.

\textbf{Tilt inside a sleeve.} For sleeve \(s\), the base weights come
from the inverse-volatility weighting of its live members. The tilt
multiplies each member's weight by a factor derived from the model's
posterior, and the sleeve's total is then restored, so the tilt changes
the composition of the sleeve and not its size.

\textbf{Shift the sleeve's weight within the bound.} The tilt may move a
sleeve's total by up to the tilt scale times the bound on the deviation.

\textbf{Hold cash.} The system's cash vehicle is a short Treasury bill
fund, which is a portfolio decision rather than an idle balance: it has
a defined net asset value, a near-zero duration, and a yield. When the
model has no edge --- and, as Chapter 20 shows, it usually does not --- the
portfolio's response is to hold more of this.

What the model may not do: set a sleeve's weight, set the total risk of
the portfolio, decide whether to hold the portfolio at all, or choose the
benchmark.

\section{How the base weights inside a sleeve are set}

Within a sleeve, the base allocation is inverse-volatility over the live
members, computed from the trailing sixty months of returns:
\begin{equation}
b_{i} = \frac{1/\hat{\sigma}_i}{\sum_{j \in s} 1/\hat{\sigma}_j},
\qquad
\hat{\sigma}_i = \sqrt{12}\, s_i,
\end{equation}
where \(s_i\) is the sample standard deviation of the fund's trailing
monthly returns, and the annualisation factor is applied for legibility
only.

Inverse-volatility is chosen over equal weight for a reason that is about
robustness rather than about performance. Equal weight over a sleeve of
eleven funds, most of which track the same index, is a perfectly
reasonable allocation and it has one very attractive property: it cannot
be much wrong. It is also, by construction, insensitive to any estimate,
which is a virtue in a layer that is downstream of a forecast the book
does not trust.

The final weights are therefore a product of two ideas with different
justifications: a stated policy\idx{Policy} for the sleeves, and a parameter-free
rule for the funds inside each sleeve, with the model's forecast permitted
to adjust the second and never the first.

\begin{keytakeaways}
\item The strategic allocation is fixed in advance, is not set by the
model, and is argued from the investor's circumstances rather than from
the data.
\item The model may tilt within a sleeve, move a sleeve's weight within a
bound, and hold cash --- and nothing else.
\item Base weights inside a sleeve are inverse-volatility over the live
members, chosen for robustness rather than for performance.
\item Everything the model does is reversible by the reader without
changing the book's argument.
\end{keytakeaways}
